\documentclass[11pt]{article}

%\usepackage[margin=0.7in, ]{geometry} 
\usepackage{amsmath,amsthm,amssymb,scrextend}
\usepackage{fancyhdr}
\usepackage{proof}
\usepackage{xargs}
\usepackage{xcolor}
\include{macros}
\pagestyle{fancy}

\begin{document}
	
Reviewer B:\\
	
p.3, l.271: Are the statement being claimed and the statement being proved not reversed? If the language determines $a_G$, then $a_G$ is invariant under the language, so the argument here does not provide a counterexample in that direction.\\
This is true, we should say that the generating function does not characterize the language, as different languages can have the same gf\\
p.6, l.602: I wonder how to prove that this fixed point is Conway. Laird (LICS 2016) proved that it is uniform, but I could not find a proof of Conway.\\ 
The proof is in Charles Grellois and Paul-Andr´e Melli`es, An infinitary model of linear logic, property 3.\\

On line 521 (p.5), "if we take on it the pointwise partial order". What do you mean by the point-wise partial order? Do you mean 
$$ \sum_\mu s_{\mu} x^\mu \le \sum_\mu t_\mu x^\mu \Leftrightarrow \forall \mu. s_\mu \le t_\mu $$?  Or, do you compare them as functions in the point-wise mannar? If the former, I am wondering if the poin-wise minimun solution computed in the proof of Theorem 6.7 is actually $a_G(1)$. If the latter, I am wondering the bijection $$ R\{\{x\}\}\{\{y\}\} \cong R\{\{x,y\}\} $$ is order-isomorphic? (On the left-hand side, are we supposed to compare the values obtained by substituting elements of $R$ for $ y $ instead of $ R\{\{x\}\}$?)\\

We indeed consider the pointwise order $ \sum_\mu s_{\mu} x^\mu \le \sum_\mu t_\mu x^\mu \Leftrightarrow \forall \mu. s_\mu \le t_\mu $. In Theorem 6.7, the minimal solution of the system is indeed $a_\gphors(1)$ for the following reason:\\
(elementary version) the coefficients $\vec a(z) \in \fps{\Rinf}{z}$ are define by fixpoint iteration as follows:
$$ \vec a^{(0)}(z) = \vec 0 \in \fps{\Rinf}{z} \quad \vec a^{(n+1)}(z) = p(\vec a^{(n)}, z) \quad \vec{a}(z)= (\sup_n \vec{a}^{(n)})(z)$$
On the other side, the minimal solution of the system $\vec x = p(\vec x , 1)$ is defined by
$$ \vec x^{(0)} = \vec 0  \quad \vec x^{(n+1)} = p(\vec x^{(n)}, 1) \quad \vec x = \sup_n \vec{x}^{(n)}$$
As evaluation in 1 is a Scott continous semiring homomorphism $\fps{\Rinf}{z} \to \Rinf$, we can prove by induction over $n$ that $\forall n \: a^{(n)}(1)$ and then $\vec{x} = a(1)$
(probably here it can be better to give a categorical statement).\\
\\2. The proof of Theorem 6.7 seems to rely on a characterzation of $$ a_{\mathcal{G}} $$ as the minimum solution of an FAS. Is it possible to decide AST sololy by the algebraicity of $$ a_{\mathcal{G}} $$ (possibly with a definifing polynomial)?\\
To our knowledge, proving 6.7 directly by algebraicity would require a deep analysis of the variable elimination procedure (Theorem 5.1) to keep track of the spurious solutions that it could introduce; as the decidability of ETOR is proved by quantifier elimination (that can be seen as a generalization of variable elimination), it looked more natural to use directly ETOR. 
It should be noted that this kind of reasonment is indeed standard in the theory of formal languages, where the celebrated Chomsky-Schutzenberger theorem is indeed proved in the same way: first one writes a system of fixpoint equations on the semiring $\mathbb{R}^{\geq 0}$ and then by variable elimination on the ring $\mathbb{R}$ one reduces on one single equation (which cannot be done on a semiring). 
While our result about termination could be proved just by considering the fixpoint equation, algebraicity of the power series has a much wider interest than AST: not only it gives an easy proof of PAST, but we can use it to decide if an algebraic PHORS has a finite value tree or two algebraic PHORS have the same tree (parellel to what is well know in the theory of context-free languages).
	
	
	
\end{document}